\documentclass[11pt]{article}
\usepackage[letterpaper,margin=1in]{geometry}
\usepackage[T1]{fontenc}
\usepackage[utf8]{inputenc}
\usepackage{lmodern,microtype}
\usepackage{amsmath,amssymb,mathtools,bm,slashed}
\usepackage{xcolor,enumitem,booktabs,array,needspace}
\usepackage[most]{tcolorbox}
\usepackage{fancyhdr,titlesec,parskip,hyperref}
\definecolor{qftblue}{HTML}{245A8D}
\definecolor{qftgreen}{HTML}{2F6B4F}
\definecolor{qftgray}{HTML}{F4F6F8}
\hypersetup{colorlinks=true,linkcolor=qftblue,urlcolor=qftblue,pdfauthor={Nima Lashkari}}
\pagestyle{fancy}\fancyhf{}
\lhead{PHYS 662: Quantum Field Theory I}\cfoot{\thepage}
\setlength{\headheight}{14pt}
\titleformat{\section}{\Large\bfseries\color{qftblue}}{\thesection.}{0.6em}{}
\titleformat{\subsection}{\large\bfseries\color{qftgreen}}{\thesubsection}{0.6em}{}
\setlist[enumerate]{leftmargin=2em,itemsep=0.25em,topsep=0.35em}
\setlist[itemize]{leftmargin=2em,itemsep=0.2em,topsep=0.35em}
\newcommand{\dd}{\mathrm d}
\newcommand{\ii}{\mathrm i}
\newcommand{\eps}{\varepsilon}
\newcommand{\tr}{\operatorname{tr}}
\newcommand{\Tr}{\operatorname{Tr}}
\newcommand{\diag}{\operatorname{diag}}
\newcommand{\Lag}{\mathcal L}
\newcommand{\D}{\mathcal D}
\newcommand{\vac}{\lvert0\rangle}
\newcommand{\vbra}{\langle0\rvert}
\newcommand{\Tord}{\mathrm T}
\newcommand{\ket}[1]{\lvert#1\rangle}
\newcommand{\bra}[1]{\langle#1\rvert}
\newtcolorbox{keyidea}[1][]{enhanced,breakable,colback=qftgray,colframe=qftblue,
 boxrule=0.7pt,arc=1.5mm,title={#1},fonttitle=\bfseries,coltitle=white}
\newtcolorbox{remarkbox}[1][]{enhanced,breakable,colback=white,colframe=qftgreen,
 boxrule=0.6pt,arc=1.5mm,title={#1},fonttitle=\bfseries,coltitle=white}
\newcommand{\module}[1]{\clearpage\section{#1}}
\newcommand{\lecturefront}[4]{%
\rhead{Lecture #1}\hypersetup{pdftitle={PHYS 662 Lecture #1: #2}}
\begin{document}\pagenumbering{gobble}
\begin{titlepage}\centering\vspace*{1.05in}
{\scshape\Large PHYS 662 --- Quantum Field Theory I\par}
\vspace{0.35in}{\Huge\bfseries\color{qftblue}Lecture #1\par}
\vspace{0.2in}{\LARGE\bfseries #2\par}
\vspace{0.2in}{\large\color{qftgreen}#3\par}
\vspace{0.6in}{\large Nima Lashkari\par}{\large Fall 2026\par}
\vfill\begin{minipage}{0.86\textwidth}\small #4\par\medskip
Course page: \url{https://nimalashkari.com/quantum-field-theory-2026-part-i/}
\end{minipage}\vspace{0.5in}\end{titlepage}
\pagenumbering{roman}\setcounter{tocdepth}{1}\tableofcontents
\vfill\begin{remarkbox}[Place in the course]
Lectures 17--18 develop the quantum-mechanical construction. Lecture 19
extends Gaussian integration and source methods to scalar fields.
Lecture 20 introduces Grassmann integration and fermionic fields.
Problem Set 6 brings these four steps together.
\end{remarkbox}\clearpage\pagenumbering{arabic}}
\rhead{Homework Set 6}
\hypersetup{pdftitle={PHYS 662 Homework Set 6: Path Integrals}}
\newlist{parts}{enumerate}{1}
\setlist[parts]{label=\textbf{(\alph*)},leftmargin=2.3em,itemsep=0.6em,topsep=0.4em}
\newcommand{\problem}[2]{\clearpage\section*{\color{qftblue}Problem #1: #2}}
\newcommand{\connection}[1]{{\small\textit{Lecture connection: #1.}\par}\medskip}
\begin{document}
\begin{center}
{\scshape\Large PHYS 662 --- Quantum Field Theory I\par}
\vspace{0.3em}{\Huge\bfseries\color{qftblue}Homework Set 6\par}
\vspace{0.3em}{\large Lectures 17--20 \quad Fall 2026\par}
\end{center}\medskip
\textbf{Name:} \rule{0.43\textwidth}{0.4pt}\hfill
\textbf{Due:} \rule{0.24\textwidth}{0.4pt}

\textbf{Instructions.} Complete all four problems, worth 25 points each.
Show the derivations, track normalization and signs, and state boundary
conditions. You may reuse a result established in an earlier part.
No exact evaluation of the interacting double-well path integral or
loop momentum integral is required.

\begin{keyidea}[Conventions]
In Problems 1--2 keep \(\hbar\) explicit and use
\[
 \langle\bm q|\bm p\rangle=e^{\ii\bm p\cdot\bm q/\hbar},\quad
 \langle\bm p|\bm p'\rangle=(2\pi\hbar)^d\delta^{(d)}(\bm p-\bm p'),
 \quad E_{\bm p}=\frac{\bm p^2}{2m}.
\]
In Problems 3--4 set \(\hbar=c=1\), use
\(\eta=\diag(1,-1,-1,-1)\), and work in \(3+1\) dimensions with \(m>0\).
Write
\[
 A=-(\Box+m^2)+\ii0,\qquad
 D_F(x-y)=\int\frac{\dd^4p}{(2\pi)^4}
 \frac{\ii e^{-\ii p\cdot(x-y)}}{p^2-m^2+\ii0}=\ii A^{-1}(x,y).
\]
All vacuum functionals have \(Z[0]=1\) unless explicitly called
unnormalized. Products such as \(JA^{-1}J\) include spacetime integrals.
Use a finite regulator before manipulating a determinant.

All Grassmann sources anticommute with fields and with each other.
Use left derivatives, with the rightmost derivative acting first, and
\[
 \int\dd\bar\theta\dd\theta\,\theta\bar\theta=1,
 \qquad \partial^L_\theta(\bar\theta\theta)=-\bar\theta.
\]
Barred and unbarred sources are differentiated independently.
\end{keyidea}

\textbf{The four problems}
\begin{enumerate}
\item The free measure, retarded propagation, and the first Born term.
\item Euclidean paths and controlled tunneling in a double well.
\item Complex scalars, Gaussian sources, and connected correlations.
\item Grassmann Gaussians, Dirac propagation, and fermionic signs.
\end{enumerate}
The double-well calculation isolates the leading nonperturbative
exponential. The scalar and fermion problems compare inverse
determinants with determinants, and symmetric pairings with
antisymmetric pairings.

\problem{1}{Free propagation and a weak localized potential}
\connection{Lecture 17, Sections 2--5; Lecture 18, Sections 1--3.}
Consider a particle of mass \(m\) in \(d\) Cartesian spatial dimensions.
Let \(T=t_f-t_i>0\), \(\epsilon=T/N\), and hold both endpoints fixed.
The potential in parts (e)--(f) is time independent.
\begin{parts}
\item \textbf{(4 points)} Starting from the momentum resolution of the
identity, derive the time-sliced configuration integral for
\(H_0=\bm p^2/(2m)\). Display all integration ranges, the number of
integrated positions and momenta, and the \(N\)-dependent prefactor.
Specify how the oscillatory Gaussian is defined.

\item \textbf{(5 points)} Shift the discretized path by its straight-line
classical solution. For the fluctuation matrix with diagonal entries
2 and nearest-neighbor entries \(-1\), prove
\(\det A_n=n+1\) by recursion. Evaluate the remaining integrals and
obtain \(K_0(\bm q_f,\bm q_i;T)\), including its phase and normalization.

\item \textbf{(4 points)} Verify the composition law and the
delta-function limit as \(T\to0^+\). Explain why the prefactor cannot
be absorbed into an arbitrary normalization depending on \(T\).
State the dimensions of the kernel.

\item \textbf{(4 points)} Define
\[
 G_R(T)=\frac{\theta(T)}{\ii\hbar}e^{-\ii H_0T/\hbar},\qquad
 \widetilde G_R(E)=\int\dd T\,e^{\ii ET/\hbar}G_R(T).
\]
Derive its unit delta-function source and show that
\(\widetilde G_R(E)=(E-H_0+\ii0)^{-1}\). Contrast this with the
equation satisfied by the evolution kernel without \(\theta(T)\).

\item \textbf{(4 points)} Expand the configuration path integral to
first order in \(V\). Derive the integral with one potential insertion
between two free kernels. Explain the time integration limits and
the factor \(-\ii/\hbar\); do not assume the Born formula in advance.

\item \textbf{(4 points)} Take
\(V(\bm q)=V_0e^{-\bm q^2/(2\ell^2)}\), with \(\ell>0\).
Remove the external free evolution phases and derive the first-order
momentum-space scattering matrix element. Evaluate the spatial
Gaussian Fourier transform and the time integral. Identify which
quantity is conserved and how the momentum-transfer scale depends
on \(\ell\). No differential cross section is requested.
\end{parts}
\begin{remarkbox}[Check]
The time integral over a static scattering insertion is
\(2\pi\hbar\delta(E_f-E_i)\). The unscattered term is
\((2\pi\hbar)^d\delta^{(d)}(\bm p_f-\bm p_i)\).
\end{remarkbox}

\problem{2}{Euclidean tunneling in a double well}
\connection{Lecture 18, Sections 4, 6--8.}
Take
\[
 H=\frac{p^2}{2m}+V(q),\qquad V(q)=\lambda(q^2-a^2)^2,
 \qquad m,\lambda,a>0.
\]
Distinguish the finite-time transition kernel from a saddle on the
infinite Euclidean time line. Work in the semiclassical regime specified
in part (d).
\begin{parts}
\item \textbf{(4 points)} Write the time-sliced Euclidean integral for
\(K_E(a,-a;\tau)=\langle a|e^{-H\tau/\hbar}|-a\rangle\), including
its measure and fixed endpoints. Derive the Euclidean Lagrangian
from the real-time action and explain the sign of its potential term.

\item \textbf{(4 points)} Derive the Euclidean Euler--Lagrange equation
and its first integral. Determine the small-oscillation frequency
\(\omega_0\) about either minimum from \(V''(\pm a)\).

\item \textbf{(6 points)} On \(-\infty<\tau<\infty\), find the
finite-action solution with \(q(-\infty)=-a\), \(q(+\infty)=a\),
and \(q'\to0\) at both ends. Display its arbitrary center \(\tau_0\).
Evaluate its action \(S_I\) explicitly in terms of \(m,\lambda,a\),
and then in terms of \(m,\omega_0,a\).
\textit{Hint:} use the conserved Euclidean energy before separating
variables. This solution reaches its minima asymptotically; it is not
an exact saddle with those endpoint values at finite times.

\item \textbf{(4 points)} Explain why the leading tunneling amplitude
contains \(e^{-S_I/\hbar}\), and give the condition that makes this
exponential small. Identify the translation zero mode of the
quadratic fluctuation operator and explain why the prefactor cannot
be inferred from \(S_I\) alone. No fluctuation determinant is required.

\item \textbf{(4 points)} Let the lowest even and odd energy eigenstates
be \(|0\rangle,|1\rangle\), with \(\Delta E=E_1-E_0>0\). In their
two-dimensional subspace define
\[
 |R\rangle=\frac{|0\rangle+|1\rangle}{\sqrt2},\qquad
 |L\rangle=\frac{|0\rangle-|1\rangle}{\sqrt2}.
\]
Compute \(\langle R|e^{-H\tau/\hbar}|L\rangle\) and
\(\langle L|e^{-H\tau/\hbar}|L\rangle\), and take their ratio.
Explain how the ratio distinguishes small and large
\(\Delta E\tau/\hbar\). These states are normalized packets in the
two-state model, not exact position eigenstates.

\item \textbf{(3 points)} You may use
\(\Delta E\simeq2\hbar\kappa e^{-S_I/\hbar}\), where \(\kappa\)
is an undetermined fluctuation factor. Determine its dimensions.
Using parity and a spectral sum, identify the large-\(\tau\) decay
rate of \(\langle0|q_E(\tau)q_E(0)|0\rangle\) when
\(\langle0|q|1\rangle\neq0\). Explain why this is a controlled
approximation rather than an exact evaluation of the double-well integral.
\end{parts}

\problem{3}{Complex scalar sources and connected functions}
\connection{Lecture 19, Sections 2--6.}
Let \(\Phi=(\phi_1+\ii\phi_2)/\sqrt2\) be a free complex scalar of
mass \(m\), with
\[
 S_0=\int\dd^4x\,
 (\partial_\mu\Phi^*\partial^\mu\Phi-m^2\Phi^*\Phi).
\]
Introduce commuting sources by \(S_J=\bar J\Phi+\Phi^*J\).
\begin{parts}
\item \textbf{(4 points)} First work at a finite Euclidean regulator
with \(n\) complex variables \(z\) and a positive Hermitian matrix \(B\).
Evaluate
\[
 \int\prod_{a=1}^n\frac{\dd\Re z_a\,\dd\Im z_a}{\pi}
                  e^{-\bar zBz}.
\]
Use the result to identify the determinant power in the unnormalized
complex-scalar vacuum functional after Feynman continuation. Compare
with one real scalar, and distinguish this answer from normalized \(Z[0]\).

\item \textbf{(5 points)} Complete the square in the regulated Minkowski
integral and derive \(Z_0[\bar J,J]\). Give the shifts of both fields,
and keep track of the factors of \(\ii\). State which boundary
condition selects the inverse of the quadratic operator.

\item \textbf{(4 points)} Obtain
\(\langle\Tord\Phi(x)\Phi^\dagger(y)\rangle\) by source derivatives.
Also evaluate \(\langle\Tord\Phi(x)\Phi(y)\rangle\). Explain the
result using both the quadratic source functional and the global
\(U(1)\) symmetry. Explain why the insertions are time ordered.

\item \textbf{(4 points)} Derive
\(\langle\Tord\Phi(x_1)\Phi(x_2)\Phi^\dagger(y_1)\Phi^\dagger(y_2)\rangle\)
by differentiating your result, explicitly listing the allowed
pairings. Apply \(\Box_x+m^2\) to the two-point function and verify
the delta-function normalization by functional integration by parts.

\item \textbf{(4 points)} For a general normalized source functional
with vanishing charged one-point functions and unbroken \(U(1)\),
define \(C=\log Z\). Expand \(Z=e^C\) through terms with two \(J\)'s
and two \(\bar J\)'s. Show that the full four-point function is the
connected four-point function plus the two products of connected
two-point functions. Use this expansion to explain why derivatives
of \(\log Z\) select connected groups at arbitrary order.

\item \textbf{(4 points)} Define \(W=-\ii\log Z\). Derive the factor
of \(\ii\) relating a total of \(r\) source derivatives of \(W\)
to the corresponding connected function, with sources coupled as
above. Apply the result at \(r=2\) and \(r=4\) in the free theory.
Explain why a nonzero full four-point function is compatible with a
vanishing connected four-point function.
\end{parts}

\problem{4}{Fermionic Gaussians, propagation, and statistics}
\connection{Lecture 20, Sections 2--8.}
Use the Grassmann conventions on the first page. Let
\(D=\ii\slashed\partial-m\), and denote its Feynman inverse by
\(R_D\), so that \(DR_D=1\). Source terms are ordered as
\(\bar\eta\psi+\bar\psi\eta\).
\begin{parts}
\item \textbf{(4 points)} Evaluate the one-pair integral
\(\int\dd\bar\theta\dd\theta\,e^{-a\bar\theta\theta}\) directly.
For two pairs and a general invertible \(2\times2\) matrix \(M\),
expand the exponential and show that its integral is \(\det M\).
Explain why the same argument produces a determinant for any number
of pairs and how this differs from a commuting complex Gaussian.

\item \textbf{(4 points)} Complete the square with odd sources and
derive the normalized Dirac functional \(Z_0[\bar\eta,\eta]\).
Track the order of barred and unbarred quantities. Identify the
unnormalized determinant before it cancels in the ratio.

\item \textbf{(5 points)} Evaluate
\[
 \left.\frac{\delta^L}{\delta\bar\eta_a(x)}
       \frac{\delta^L}{\delta\eta_b(y)}Z_0\right|_0
 \quad\hbox{and the expression with the derivatives reversed}.
\]
Relate them to \(S_{F,ab}(x-y)=\langle\Tord\psi_a(x)\bar\psi_b(y)\rangle\).
Derive its momentum-space form and check
\(D_xS_F(x-y)=\ii\delta^{(4)}(x-y)I_4\). Verify this source term
also using the fermionic Schwinger--Dyson identity.

\item \textbf{(4 points)} Derive the two Wick pairings of
\(\langle\Tord\psi_a(x_1)\bar\psi_b(y_1)
                 \psi_c(x_2)\bar\psi_d(y_2)\rangle\),
including their relative sign. Contract indices to obtain the
connected two-point function of
\(j^\mu=\bar\psi\gamma^\mu\psi\). Explain the closed-loop minus sign.

\item \textbf{(4 points)} For one mode with \(H=\omega c^\dagger c\),
evaluate the thermal trace using both the empty/occupied basis and
\[
 \Tr O=\int\dd\bar\xi\dd\xi\,e^{-\bar\xi\xi}
                        \langle-\bar\xi|O|\xi\rangle.
\]
Use the coherent kernel from Lecture 20 to check agreement. Explain
how the minus sign enforces antiperiodic paths and derive their
allowed Euclidean frequencies. What trace results if this sign is omitted?

\item \textbf{(4 points)} Couple to a prescribed background through
\(\mathcal L_{\rm int}=-e\bar\psi\gamma^\mu\psi A_\mu\).
Expand \(\log[\det(D-e\slashed A)/\det D]\) to second order in
\(A\). Identify the cyclic spinor trace and explain its agreement
with part (d). State the vertex factor and the extra sign for each
closed fermion loop. Do not evaluate the divergent momentum integral.
\end{parts}
\end{document}
